\begin{abstract}
Ushbu tadqiqotda "AI Agents for Scientific Discovery" mavzusi bo'yicha jahon ilmiy adabiyotlari, jumladan arXiv, CrossRef va PubMed xalqaro bazalaridagi zamonaviy ilmiy ishlar tizimli tahlil qilindi. Tadqiqot doirasida ilgari surilgan asosiy metodologiyalar, arxitekturalar va empirik natijalar qiyosiy matritsa asosida umumlashtirildi. Mavjud yondashuvlarning kuchli va zaif tomonlari baholanib, kelgusidagi istiqbolli tadqiqot yo'nalishlari belgilab berildi.

\textbf{Kalit so'zlar:} AI Agents for Scientific Discovery, adabiyotlar tahlili, tizimli sharh, metodologiya, sun'iy intellekt.
\end{abstract}
